\section{Unit testy}

\subsection{Charakteristika unit testov}
Unit testy sú automatizované testy, ktoré overujú správnosť izolovaných častí kódu — konkrétnych funkcií alebo tried — bez závislosti na externých systémoch, ako sú Firebase, Google Maps alebo používateľské rozhranie. Každý test funguje na overenom princípe:

\begin{enumerate}[noitemsep]
      \item \textbf{Arrange:} Príprava vstupných dát a stavu.
      \item \textbf{Act:} Vykonanie testovanej funkcie.
      \item \textbf{Assert:} Porovnanie výsledku s očakávanou hodnotou.
\end{enumerate}

  Ak sa výsledok zhoduje s očakávaním, test je vyhodnotený ako \textbf{PASSED}. V opačnom prípade zlyhá (\textbf{FAILED}) a poskytne informáciu o mieste a príčine chyby.

  \subsection{Umiestnenie a spustenie testov}
  Všetky unit testy sú umiestnené v module \texttt{app/src/test/} a spúšťajú
  sa príkazom \texttt{./gradlew testDebugUnitTest}. Po vykonaní systém
  vygeneruje HTML správu s detailnými výsledkami.

  \subsection{Pomocná funkcia testEvent}
  Pre optimalizáciu kódu a elimináciu redundancie pri vytváraní testovacích objektov bola implementovaná pomocná funkcia \texttt{testEvent()}. Táto funkcia obsahuje predvolené hodnoty pre všetky povinné parametre triedy        
  \texttt{Event}, čo umožňuje v jednotlivých testoch definovať len tie atribúty, ktoré sú pre daný scenár relevantné.

  \subsection{Popis testovacích skupín}

  \subsubsection{Skupina 1: NoShowCountTest}
  Zameriava sa na funkciu \texttt{UserProfile.recentNoShowCount()}, ktorá počíta neúčasti používateľa za posledných 30 dní. Táto logika je kritická pre správne fungovanie penalizačného systému (ban po 3 neúčastiach).

  \begin{table}[H]
      \centering                                                                                                                                                                                                                   
      \caption{Prehľad testov pre NoShowCount}
      \begin{tabular}{|p{4.5cm}|p{4cm}|c|p{3.5cm}|}
          \hline                                                                                                                                                                                                                   
          \textbf{Názov testu} & \textbf{Vstup (noShowTimestamps)} & \textbf{Očak.} & \textbf{Účel} \\ \hline                                                                                                                      
          no timestamps returns zero & prázdny zoznam & 0 & Overenie nového užívateľa. \\ \hline                                                                                                                                   
          two recent no-shows & \texttt{listOf(now-1s, now-2s)} & 2 & Základné počítanie. \\ \hline                                                                                                                                
          older than 30 days ignored & \texttt{listOf(now - 31 dní)} & 0 & Overenie časového limitu. \\ \hline                                                                                                                     
          mix of old and recent & \texttt{listOf(31d, 1s, 2s)} & 2 & Filtrovanie starých záznamov. \\ \hline                                                                                                                       
      \end{tabular}
  \end{table}

  \subsubsection{Skupina 2: CancelDeadlineTest}
  Overuje funkciu \texttt{Event.isPastCancelDeadline()}, ktorá určuje, či používateľ ešte môže zrušiť svoju účasť bez penalizácie.

  \begin{table}[H]
      \centering                                                                                                                                                                                                                   
      \caption{Prehľad testov pre zrušenie účasti}
      \begin{tabular}{|p{4.5cm}|p{4cm}|c|p{3.5cm}|}
          \hline                                                                                                                                                                                                                   
          \textbf{Scenár} & \textbf{Parametre} & \textbf{Výsledok} & \textbf{Logika} \\ \hline                                                                                                                                     
          checkInType "none" & akýkoľvek čas & \texttt{false} & Bez check-inu niet postihu. \\ \hline                                                                                                                              
          deadline set to 0 & akýkoľvek čas & \texttt{false} & Možnosť zrušiť kedykoľvek. \\ \hline                                                                                                                                
          past deadline & Event o 1h, deadline 3h & \texttt{true} & Odhlásenie po termíne. \\ \hline                                                                                                                               
          before deadline & Event o 10h, deadline 2h & \texttt{false} & Bezpečné odhlásenie. \\ \hline                                                                                                                             
      \end{tabular}
  \end{table}

  \subsubsection{Skupina 3: EventFilterTest}
  Testuje \texttt{EventFilterService}, ktorý filtruje zoznam podujatí na základe užívateľských preferencií. Pre účely testu sú definované 4 vzorové podujatia (Jazz, Futbal, Yoga, Tech).

  \begin{itemize}[noitemsep]
      \item \textbf{Filter by name:} Overuje, že vyhľadávanie je \textit{case-insensitive} (napr.\ „jazz" nájde „Jazz Evening").
      \item \textbf{Filter by category:} Overuje správnosť zaradenia do kategórií (napr.\ „Sport" vráti Futbal aj Yogu).
      \item \textbf{Filter by maxPrice:} Testuje numerické porovnávanie cien (napr.\ do 5\,€ neukáže konferenciu za 20\,€).
  \end{itemize}

  \subsubsection{Skupina 4: ActiveEventsTest}
  Overuje logiku filtrovania aktívnych podujatí pre zobrazenie na mape.
  \begin{itemize}[noitemsep]
      \item \textbf{Past event:} Podujatie, ktoré už skončilo, sa nesmie vrátiť.
      \item \textbf{Future event:} Budúce podujatia sa musia vrátiť v zozname.
      \item \textbf{No date event:} Podujatia bez dátumu sú považované za trvalo aktívne.
  \end{itemize}

  \subsubsection{Skupina 5: RecommendationEngineTest}
  Testuje komplexný algoritmus triedy \texttt{RecommendationEngine}, ktorý zoraďuje podujatia podľa relevancie.

  \begin{table}[H]
      \centering                                                                                                                                                                                                                   
      \caption{Logika odporúčacieho algoritmu}
      \begin{tabular}{|p{5cm}|p{8cm}|}
          \hline                                                                                                                                                                                                                   
          \textbf{Testovaný faktor} & \textbf{Očakávané správanie} \\ \hline                                                                                                                                                       
          \textbf{Vzdialenosť} & Bližšie podujatie musí mať vyšší rank ako vzdialené. \\ \hline                                                                                                                                    
          \textbf{Navštívené} & Podujatia v \texttt{visitedEventIds} sú z výsledkov vylúčené. \\ \hline                                                                                                                            
          \textbf{Vlastné} & Podujatia vytvorené aktuálnym používateľom sa mu neodporúčajú. \\ \hline                                                                                                                              
          \textbf{Minulosť} & Ukončené podujatia sa do odporúčaní nedostanú. \\ \hline                                                                                                                                             
      \end{tabular}
  \end{table}

  \subsection{Súhrnný výsledok testovania}
  Celkovo bolo vykonaných \textbf{24 unit testov}.

  \begin{quote}
      \texttt{BUILD SUCCESSFUL} \\                                                                                                                                                                                                 
      \texttt{24 tests completed, 0 failed}
  \end{quote}

  Úspešné absolvovanie všetkých testov potvrdilo stabilitu a správnosť
  kľúčovej biznis logiky aplikácie Joinly, najmä v oblastiach penalizácie
  používateľov, filtrovania obsahu a personalizovaných odporúčaní